\subsection{\bt{Maintenance cost}{维护代价}}
\label{sec:cost}
\bi{\emph{Controlled benchmark (Q4).} The benchmark admits $k$ structured definitions, $k\in\{1,2,4,6\}$, from each of four families (sales, returns, return ratios, catalog), 4 to 19 in all, with $k=1$ retaining sharing between the sales and return-ratio families. After each of four updates, three benign appends and the status-change rows, eight scripted agents use every definition in shuffled order, with staggered or simultaneous arrivals, and database time sums the operations during these uses. Besides the full recheck per definition and \system, we run a variant that rechecks changed tables only, reusing no check results, and the query cache, which caches the full recheck's check queries by normalized SQL text and table versions, with no knowledge of condition types.}
{\emph{受控实验（Q4）。}基准从四个族（销售、退货、退货比率、目录）各准入最多 $k\in\{1,2,4,6\}$ 个结构化定义，共 4 到 19 个，$k=1$ 仍保留销售与退货比率族之间的共享。四次更新（三次正常追加，以及退货改为状态流水）之后，8 个脚本智能体各按打乱的顺序使用所有定义，错峰或同时到达，数据库时间累计这些使用期间的操作。除整定义重查与 \system\ 外，还运行只重查变化的表的变体（不复用检查结果）与查询缓存（按规范化 SQL 文本与表版本缓存整定义重查的检查查询，不了解条件类型）。}

\begin{figure}[t]
\centering
\begin{tikzpicture}
\begin{groupplot}[group style={group size=2 by 1, horizontal sep=6mm, y descriptions at=edge left},
  mavra, width=0.6\linewidth, height=42mm, xmin=2, xmax=21, ymin=0, ymax=125,
  xtick={4,8,15,19}, ytick={0,25,50,75,100}, yticklabel={\pgfmathprintnumber{\tick}\%},
  xlabel={\bt{Definitions}{定义数}}, title style={font=\footnotesize, yshift=-1ex}]
\nextgroupplot[title={\bt{Staggered}{错峰}}, ylabel={\bt{DB time vs.\ definition-level}{相对定义级的 DB 时间}},
  legend to name=costlegend, legend columns=\ifbilingual 2\else 4\fi, legend style={font=\scriptsize, draw=none, column sep=3pt}]
% generated by tools/review-results.py
\addplot[costScope] coordinates {(4,94.4) (8,82.9) (15,84.3) (19,86.8)};
\addplot[costCache] coordinates {(4,56.8) (8,30.3) (15,18.2) (19,16.7)};
\addplot[costCond] coordinates {(4,42.5) (8,24.8) (15,17.0) (19,14.0)};

\addplot[mDef, line width=0.9pt, densely dotted, domain=2:21, samples=2] {100};
\addlegendentry{\bt{Scope-only}{仅范围}}
\addlegendentry{\bt{Definition + cache}{定义级 + 缓存}}
\addlegendentry{\system}
\addlegendentry{\bt{Definition-level}{定义级}}
\nextgroupplot[title={\bt{Simultaneous}{同时到达}}]
% generated by tools/review-results.py
\addplot[costScope] coordinates {(4,86.1) (8,90.0) (15,91.5) (19,119.3)};
\addplot[costCache] coordinates {(4,71.3) (8,84.3) (15,59.3) (19,50.0)};
\addplot[costCond] coordinates {(4,78.2) (8,84.1) (15,61.1) (19,53.1)};

\addplot[mDef, line width=0.9pt, densely dotted, domain=2:21, samples=2] {100};
\end{groupplot}
\node[anchor=south] at ($(group c1r1.north)!0.5!(group c2r1.north)+(0,4mm)$) {\pgfplotslegendfromname{costlegend}};
\end{tikzpicture}
\caption{\bt{Maintenance database time relative to definition-level revalidation (100\%), 1\,M sales rows, eight agents. The generic cache keys check results by SQL text and table versions; \system\ keys verdicts by condition identity and versions.}{相对定义级重验证（100\%）的维护数据库时间，100 万行销售，8 个智能体。通用缓存按 SQL 文本与表版本索引检查结果；\system\ 按条件身份与版本索引结论。}}
\label{fig:cost}
\Description{Two line charts of database time relative to definition-level revalidation over 4 to 19 definitions. Staggered: MAVRA and the generic cache fall from about 43--57\% to 14--17\%; scope-only stays at 83--94\%. Simultaneous: MAVRA and the cache lie between 50\% and 84\%; scope-only between 86\% and 119\%.}
\end{figure}


\bi{Figure~\ref{fig:cost} shows that check-result reuse, not scope, provides the savings, and that it does not require typed conditions. Under staggered arrivals, the full recheck's time grows with the number of definitions (\CbStagDefOne\ to \CbStagDefSix\,s), whereas \system\ and the query cache stay at \CbCondStagMin--\CbCondStagMax\ and \CbCacheStagMin--\CbCacheStagMax\,s; with 19 definitions they need \CbStagCondSix\% and \CbStagCacheSix\% of the full recheck's time, and rechecking changed tables only needs 83--94\%. The one visible gap appears with the least sharing (\CbStagCondOne\% vs.\ \CbStagCacheOne\%), where \system's canonical form recognizes a condition just checked with its key columns in another order. Under simultaneous arrivals, request coalescing already removes much duplicate work (50--84\%). Without shared conditions, all methods cost \CbZeroMin--\CbZeroMax\% of the full recheck's time: the savings come from sharing alone. All cells record zero stale uses, false invalidations, and unavailable uses.}
{图~\ref{fig:cost} 表明，节省来自检查结果复用而不是缩小检查范围，而且不需要类型化条件。错峰到达时，整定义重查的耗时随定义数增长（\CbStagDefOne\ 到 \CbStagDefSix\ 秒），\system\ 与查询缓存则保持在 \CbCondStagMin--\CbCondStagMax\ 与 \CbCacheStagMin--\CbCacheStagMax\ 秒；19 个定义时二者分别只需整定义重查的 \CbStagCondSix\% 与 \CbStagCacheSix\%，只重查变化的表需要 83--94\%。唯一明显的差距出现在共享最少时（\CbStagCondOne\% 对 \CbStagCacheOne\%）：\system\ 的规范形式能识别以另一种键列顺序刚检查过的同一条件。同时到达时，请求合并已去掉大量重复工作（50--84\%）。没有共享条件时，各方法都需要整定义重查的 \CbZeroMin--\CbZeroMax\%：节省只来自共享。所有配置都是零过期使用、零误失效、零不可用。}

\bi{\emph{Scale and concurrency.} The ratios persist with data size and agent count: at 4 million rows \system\ and the cache need \CbFourStagCond\% and \CbFourStagCache\% of the full recheck's time under staggered arrivals, with 32 agents \CbAgentsStagCond\% and \CbAgentsStagCache\%, and from 1 to 16 million rows \system's share stays at \EoneRatioSixteen--\EoneRatioFour\% (\EoneCondSixteen\ versus \EoneDefSixteen\,s at 16 million rows).}
{\emph{规模与并发。}这些比例在数据量与智能体数增加后保持不变：400 万行时，错峰到达下 \system\ 与查询缓存分别需要整定义重查的 \CbFourStagCond\% 与 \CbFourStagCache\%；32 个智能体时为 \CbAgentsStagCond\% 与 \CbAgentsStagCache\%；从 100 万到 1600 万行，\system\ 的比例稳定在 \EoneRatioSixteen--\EoneRatioFour\%（1600 万行时 \EoneCondSixteen\ 秒对 \EoneDefSixteen\ 秒）。}

